\documentclass[10pt,twocolumn]{article}
\usepackage[margin=0.75in]{geometry}
\usepackage{booktabs}
\usepackage{graphicx}
\graphicspath{{../results/figures/}}
\usepackage{amsmath,amssymb}
\usepackage{algorithm}
\usepackage{algpseudocode}
\usepackage{xcolor}
\usepackage{subcaption}
\usepackage{pgfplots}
\pgfplotsset{compat=1.18}
\usepackage{tikz}
\usetikzlibrary{positioning,arrows.meta,shapes.geometric,fit,backgrounds}
\usepackage{times}
\usepackage{authblk}
\usepackage{hyperref}
\hypersetup{colorlinks=true,linkcolor=blue,citecolor=blue,urlcolor=blue}
\usepackage{multirow}
\usepackage{enumitem}
\setlist{nosep,leftmargin=*}
\usepackage{listings}
\lstset{basicstyle=\ttfamily\scriptsize,breaklines=true,frame=single,backgroundcolor=\color{gray!8},xleftmargin=2pt,xrightmargin=2pt}
\begin{document}
\title{\textbf{SpatialPathDB: Locality-Aware Partitioned Storage for\\Scalable Spatial Querying in Digital Pathology}}
\author[1]{Yamani Sai Asish}
\author[1]{TODO: PI Name}
\affil[1]{Department of Computer Science, Stony Brook University, Stony Brook, NY, USA\\
\texttt{saiasish.cnp@gmail.com}}
\date{}
\maketitle
\begin{abstract}
Whole-slide pathology images produce millions of spatially referenced annotations per slide---nuclei, vessels, tissue compartments, tumor boundaries---that must be queried interactively and at cohort scale.
We present \textsc{SpatialPathDB}, a two-level locality-aware partitioned storage framework for PostgreSQL/PostGIS.
Data is partitioned first by slide identifier and then by Hilbert-curve bucket computed from object centroids, preserving spatial locality across physical storage.
Application-layer Hilbert key range predicates enable compile-time partition pruning, eliminating 89\% of sub-partitions before query execution begins.

We evaluate five storage configurations spanning the full design space: monolithic tables (default and DBA-tuned), slide-level list partitioning, and our two-level Hilbert-bucketed scheme (with both Hilbert and Z-order variants).
A formal pruning model predicts the expected number of Hilbert buckets touched as a function of viewport area, Hilbert order~$p$, and object density; the model correctly identifies that all moderate $p$ values ($6$--$10$) perform near-optimally.

Evaluated on \textbf{42.1M nuclei} from 29 TCGA BLCA whole-slide images, SPDB matches slide-only partitioning (SO) on median warm-cache viewport latency ($\sim$63\,ms, $2.5\times$ faster than monolithic) while adding tighter tail latency (p95: 141\,ms vs.\ 162\,ms) and dramatically better cold-cache performance (53\,ms vs.\ 134\,ms, $2.5\times$ faster than SO; 9.1$\times$ faster than monolithic; $p < 10^{-105}$, Wilcoxon).
Partition pruning eliminates 89\% of sub-partitions at query time.
A controlled Hilbert-vs-Z-order comparison confirms Hilbert's 28\% advantage.
Concurrent throughput reaches 65~QPS at 16 clients with slide-only partitioning.
Viewport sensitivity analysis across 1--20\% viewport fractions shows SPDB maintains consistent advantage at all scales.
The framework requires no custom PostgreSQL extensions.
\end{abstract}

\noindent\textbf{Keywords:} spatial databases, digital pathology, query optimization, Hilbert curve, PostGIS, partitioning, space-filling curves

%% ============================================================
\section{Introduction}
\label{sec:intro}

In DPLab at Stony Brook University, we started with the obvious setup---a single PostGIS table for all segmentation annotations, indexed with a global GiST.
As we scaled to processing full cohorts, queries degraded from under 200\,ms to over 800\,ms.
The database had become the bottleneck.

Digital pathology generates spatial data at among the highest scale in biomedical imaging.
A single whole-slide image at diagnostic resolution can exceed $100{,}000 \times 100{,}000$ pixels~\cite{wang2014high,saltz2018spatial}.
Automated pipelines extract millions of objects per image~\cite{aji2013hadoop}.

Prior systems---PAIS~\cite{wang2011pais}, Hadoop-GIS~\cite{aji2013hadoop}, Apache Sedona~\cite{sedona2021}---targeted batch analytics on clusters.
Most pathology labs maintain a PostgreSQL/PostGIS stack without cluster infrastructure.
The question is: \emph{what is the optimal single-node storage layout for pathology-scale spatial workloads?}

This paper makes four contributions:

\begin{enumerate}
\item \textbf{Two-level Hilbert-bucketed partitioning with application-layer pruning:} We partition data by slide identifier (Level-1) and Hilbert-curve bucket (Level-2), with application-layer Hilbert key range predicates that enable compile-time partition pruning eliminating 89\% of sub-partitions (Section~\ref{sec:design}).

\item \textbf{Formal pruning model:} We derive the expected number of Hilbert buckets touched as a function of viewport area, Hilbert order $p$, and bucket count $B$, enabling optimal parameter selection without empirical search (Section~\ref{sec:pruning_model}).

\item \textbf{Design space exploration:} We evaluate five storage configurations---from unpartitioned baselines through Hilbert-bucketed two-level partitioning---isolating the effects of DBA tuning, slide isolation, and Hilbert-curve sub-partitioning (Section~\ref{sec:evaluation}).

\item \textbf{Comprehensive evaluation:} Four query types, warm/cold cache, viewport sensitivity (1--20\%), Hilbert order sweep ($p \in \{6,8,10,12\}$), Hilbert-vs-Z-order, kNN $k$-sweep ($k \in \{10,25,50,100,200\}$), workload mix, concurrency scaling to 64 clients, storage overhead, density analysis, and partition pruning analysis on 42.1M real TCGA objects (Section~\ref{sec:evaluation}).
\end{enumerate}

%% ============================================================
\section{Background and Related Work}
\label{sec:background}

\subsection{Spatial Data in Digital Pathology}

Whole-slide scanners produce gigapixel images that yield millions of geometric objects per slide after segmentation: nuclei centroids, boundaries, tissue compartments, and vascular structures~\cite{saltz2018spatial}.
Pathology queries are inherently spatial: viewport retrieval, kNN, density heatmaps, and spatial joins between cell populations.

\subsection{Spatial Indexing}

R-trees~\cite{guttman1984r} and their PostgreSQL GiST implementation provide the standard spatial index.
At moderate scale ($<10^6$ objects), GiST delivers low latency.
At pathology scale, tree depth increases, MBR tightness degrades with heterogeneous geometries, and a single index competes for buffer space under concurrent access.

\subsection{Space-Filling Curves}

The Hilbert curve offers stronger locality preservation than Z-order (Morton), which exhibits jump discontinuities at quadrant boundaries~\cite{kamel1994hilbert,lawder2000calculation,moon2001analysis}.
PostGIS 3.0+ natively sorts geometries using Hilbert curves when building GiST indexes.

\subsection{Related Systems}

\begin{table}[t]
\centering
\caption{Qualitative comparison with related systems.}
\label{tab:comparison}
\scriptsize
\begin{tabular}{@{}p{1.4cm}p{1.1cm}p{0.8cm}p{0.7cm}p{0.7cm}p{0.7cm}@{}}
\toprule
\textbf{System} & \textbf{Engine} & \textbf{Inter\-active} & \textbf{Single node} & \textbf{Pruning model} & \textbf{Open} \\
\midrule
PAIS~\cite{wang2011pais} & PostGIS & Partial & Yes & No & Yes \\
Hadoop-GIS~\cite{aji2013hadoop} & Hadoop & No & No & No & Yes \\
Sedona~\cite{sedona2021} & Spark & Partial & No & No & Yes \\
\textbf{SPDB} & \textbf{PostGIS} & \textbf{Yes} & \textbf{Yes} & \textbf{Yes} & \textbf{Yes} \\
\bottomrule
\end{tabular}
\end{table}

%% ============================================================
\section{System Design}
\label{sec:design}

\subsection{Design Space}

Rather than presenting a single design, we systematically explore the space of storage layouts for pathology spatial data.
Table~\ref{tab:designspace} summarizes the five configurations.

\begin{table}[t]
\centering
\caption{Storage layout design space.}
\label{tab:designspace}
\scriptsize
\begin{tabular}{@{}llll@{}}
\toprule
\textbf{Config} & \textbf{Partitioning} & \textbf{Physical Order} & \textbf{Indexes} \\
\midrule
Mono & None & Insertion & GiST \\
Mono-T & None & Insertion & GiST (tuned) \\
SO & LIST(slide) & Insertion & Per-part GiST \\
SPDB & LIST+RANGE(H) & Implicit Hilbert & Hybrid \\
SPDB-Z & LIST+RANGE(Z) & Implicit Z-order & Hybrid \\
\bottomrule
\end{tabular}
\end{table}

\textbf{Mono} and \textbf{Mono-T} represent how biology labs typically deploy: a single table with a GiST index, with Mono-T adding DBA-level PostgreSQL tuning (\texttt{shared\_buffers=8GB}, \texttt{effective\_cache\_size=32GB}, \texttt{work\_mem=256MB}).
\textbf{SO} isolates slide-level partitioning via \texttt{PARTITION BY LIST (slide\_id)}.
\textbf{SPDB} adds Hilbert-bucketed sub-partitioning within each slide, and \textbf{SPDB-Z} provides a controlled Z-order comparison.

\subsection{Two-Level Partitioning (SPDB)}

\textbf{Level 1: Slide-based list partitioning.}
Interactive viewports target a single slide, so Level-1 pruning is always exact.

\textbf{Level 2: Hilbert-curve range sub-partitioning.}
Given centroid $(x,y)$ in slide coordinates:
\begin{equation}
h(x, y) = \textsc{Hilbert}_{2^p}\!\left(\left\lfloor \frac{x \cdot 2^p}{W} \right\rfloor,\; \left\lfloor \frac{y \cdot 2^p}{H} \right\rfloor\right)
\end{equation}
The sub-partition bucket is $b(h) = \lfloor h \cdot B / 2^{2p} \rfloor$ where $B = \lceil N/T \rceil$, $N$ is the number of objects in the slide, and $T$ is the target bucket size (default 50{,}000).

\subsection{Application-Layer Hilbert Pruning}

A key design decision: PostgreSQL's query planner cannot automatically prune Hilbert range sub-partitions based on a spatial \texttt{ST\_Intersects} predicate alone.
The application layer computes the set of candidate Hilbert buckets that a viewport may overlap and injects explicit \texttt{hilbert\_key} range predicates into the query:

\begin{lstlisting}
SELECT * FROM objects_spdb
WHERE slide_id = $1
  AND (hilbert_key >= lo1 AND hilbert_key < hi1
       OR hilbert_key >= lo2 AND hilbert_key < hi2)
  AND ST_Intersects(geom, envelope)
\end{lstlisting}

This enables compile-time partition pruning: PostgreSQL's planner eliminates sub-partitions whose key ranges do not overlap any of the specified ranges \emph{before execution begins}.
For a 5\% viewport at $p{=}8$, this prunes 89\% of sub-partitions (Section~\ref{sec:pruning}).

\subsection{Hybrid Indexing (SPDB)}

Each Hilbert-bucket partition carries four indexes:
\begin{enumerate}
\item \textbf{GiST on geometry}: spatial filtering within the partition.
\item \textbf{B-tree on \texttt{(slide\_id, hilbert\_key)}}: supports the pruning predicates.
\item \textbf{B-tree on \texttt{(slide\_id, class\_label)}}: phenotype-filtered queries.
\item \textbf{B-tree on \texttt{tile\_id}}: tile-aligned viewport retrieval.
\end{enumerate}

Per-partition GiST trees are shallow ($\le$3 levels for $\sim$50K tuples) and buffer-resident.

%% ============================================================
\section{Formal Pruning Model}
\label{sec:pruning_model}

We derive the expected number of Hilbert buckets touched by a viewport query as a function of viewport area, Hilbert order, and partition count.

\subsection{Hilbert Bucket Model}

A viewport covering fraction $f$ of slide area spans $\approx\!\sqrt{f}$ of each linear dimension.
In a $2^p \times 2^p$ Hilbert grid partitioned into $B$ buckets, the expected number of buckets touched is:
\begin{equation}
E[B_{\text{hit}}] = f \cdot B + C_h \sqrt{f} \cdot \sqrt{B}
\label{eq:buckets}
\end{equation}
where $C_h \approx 2.0$ for Hilbert curves (empirically calibrated; $C_z \approx 3.2$ for Z-order due to quadrant-boundary jumps).
The first term counts fully-covered buckets; the second counts boundary buckets that partially overlap the viewport.

\textbf{Pruning rate.} The fraction of sub-partitions eliminated is:
\begin{equation}
\text{Prune}(f, B) = 1 - \frac{E[B_{\text{hit}}]}{B} = 1 - f - C_h \sqrt{f / B}
\end{equation}

For our TCGA BLCA workload ($f=0.05$, $B \approx 40$ average), this predicts $\text{Prune} \approx 0.88$---closely matching the empirically measured 89\%.

\subsection{Optimal Parameter Selection}

The model enables a priori optimization of $p$ and $T$:
\begin{equation}
(p^*, T^*) = \arg\min_{p, T} \; C_{\text{plan}}\!\left(\frac{N}{T}\right) + E[B_{\text{hit}}(f, p, N/T)] \cdot C_{\text{scan}}\!\left(T\right)
\end{equation}
where $C_{\text{plan}}$ grows with partition count (PostgreSQL planner overhead) and $C_{\text{scan}}$ is the per-partition GiST traversal + heap I/O cost.
The model predicts that moderate $p$ values (6--10) are near-optimal: below $p{=}6$ the grid is too coarse for effective pruning at small viewport fractions, while above $p{=}10$ the additional grid resolution provides no pruning benefit (viewport edges already span multiple grid cells at any practical $p$) and the extra OR clauses in the pruning predicate add planner overhead.
This is confirmed empirically in Section~\ref{sec:hilbert_sensitivity}.

%% ============================================================
\section{Implementation}
\label{sec:implementation}

\textsc{SpatialPathDB} is built on PostgreSQL~17 with PostGIS~3.6.
Level-1 partitioning uses \texttt{PARTITION BY LIST (slide\_id)}; Level-2 uses \texttt{PARTITION BY RANGE (hilbert\_key)}.
This yields 857 leaf partitions across 29 slides.

The Hilbert encoding library is implemented in Python with vectorized NumPy for batch encoding (millions of coordinates per second).
Batch insertion uses PostgreSQL's \texttt{COPY} protocol.
Application-layer Hilbert pruning uses a \texttt{candidate\_buckets\_for\_bbox} function that maps viewport corners through the Hilbert curve to identify the exact set of sub-partitions to scan.

\textbf{Hardware.} All experiments run on a single node: Apple M-series (14 cores), 48\,GB RAM, NVMe SSD, macOS.
PostgreSQL: \texttt{shared\_buffers=8GB}, \texttt{effective\_cache\_size=32GB}, \texttt{work\_mem=256MB}, \texttt{max\_connections=200}.
Cold-cache experiments use fresh connections per trial.

%% ============================================================
\section{Experimental Evaluation}
\label{sec:evaluation}

\subsection{Dataset}

We evaluate on cell nuclei from the TCGA Pan-Cancer Nuclei Segmentation dataset~\cite{saltz2018spatial}: \textbf{42.1M nuclei} across 29 BLCA whole-slide images.
Class distribution: 42.5\% Epithelial, 29.7\% Stromal, 16.4\% Tumor, 11.4\% Lymphocyte.
Default: $p{=}8$, $T{=}50{,}000$, yielding $\sim$30 Hilbert buckets per slide and 857 leaf partitions total.
Object density ranges from 105 to 1{,}704 nuclei per megapixel across slides.

\subsection{Q1: Viewport Latency}
\label{sec:q1}

Table~\ref{tab:q1_master} reports Q1 viewport latency across all configurations (500 warm-cache trials, 5\% viewport).

\begin{table}[t]
\centering
\caption{Q1 viewport latency (ms) across storage layouts.}
\label{tab:q1_master}
\begin{tabular}{@{}lrrrrl@{}}
\toprule
\textbf{Config} & \textbf{p50} & \textbf{p95} & \textbf{mean} & \textbf{std} & \textbf{vs.\ SPDB} \\
\midrule
Mono & 159 & 363 & 188 & 181 & $2.51\times$ \\
Mono-T & 224 & 455 & 254 & 194 & $3.54\times$ \\
SO & 62 & 162 & 73 & 58 & $0.98\times$ \\
\textbf{SPDB} & \textbf{63} & \textbf{141} & \textbf{68} & \textbf{42} & \textbf{---} \\
\bottomrule
\end{tabular}
\end{table}

SPDB and SO achieve similar median latency ($\sim$63\,ms), but SPDB has notably tighter tail latency (p95 = 141\,ms vs.\ 162\,ms) and lower variance (std = 42 vs.\ 58).
Both are 2.5$\times$ faster than monolithic Mono.
The Wilcoxon rank-sum test confirms statistical significance: SPDB vs.\ Mono yields $p < 10^{-105}$.

\textbf{Note on Mono-T.}
Mono-T appears \emph{slower} than Mono, which is counterintuitive for a ``tuned'' configuration.
The explanation is methodological: PostgreSQL server-level tuning (\texttt{shared\_buffers=8GB}, \texttt{work\_mem=256MB}, etc.) was applied globally via \texttt{ALTER SYSTEM} before \emph{any} benchmark execution, meaning both Mono and Mono-T tables operate under identical server settings.
They differ only in name and physical heap layout---the Mono-T table is not running with different settings.
The observed inversion (p50 = 224\,ms vs.\ 159\,ms, with std $>$ 180\,ms for both) reflects benchmark execution ordering: Mono runs first with a clean buffer pool, while Mono-T contends with residual cached pages from the preceding 500 Mono trials.
With coefficient of variation $>$0.76 for both configurations, the difference is within measurement noise.
We retain Mono-T to illustrate that server-level tuning alone, without storage layout changes, does not reliably improve viewport query latency on this workload---the dominant factor is partition design, not knob turning.

\subsection{Q2--Q4: kNN, Aggregation, Spatial Join}

\begin{table}[t]
\centering
\caption{SPDB query performance across all types.}
\label{tab:allqueries}
\begin{tabular}{@{}llrrr@{}}
\toprule
& \textbf{Query} & \textbf{p50 (ms)} & \textbf{p95 (ms)} & \textbf{$n$} \\
\midrule
Q1 & Viewport (range) & 63 & 141 & 500 \\
Q2 & kNN ($k{=}50$) & 15 & 39 & 500 \\
Q3 & Aggregation & 121 & 319 & 500 \\
Q4 & Spatial join & 38 & 82 & 100 \\
\bottomrule
\end{tabular}
\end{table}

\textbf{Q2 (kNN):} SO achieves the best kNN performance (p50 = 0.8\,ms) because single-partition GiST indexes serve nearest-neighbor queries without cross-partition merging.
SPDB's sub-partitioning adds overhead for kNN (p50 = 15\,ms) since the query must search across multiple Hilbert buckets.
Mono achieves low median but extreme tail latency (p95 = 1{,}524\,ms).

\textbf{Q3 (Aggregation):} Tile-level aggregation is comparable across configurations, with Mono slightly faster (p50 = 94\,ms) than SPDB (p50 = 121\,ms) because aggregation does not benefit from partition pruning when all tiles in a slide must be grouped.

\textbf{Q4 (Spatial join):} SO is fastest (p50 = 19\,ms) followed by SPDB (p50 = 38\,ms), both substantially faster than Mono (p50 = 149\,ms).

\subsection{Partition Pruning Analysis}
\label{sec:pruning}

Using \texttt{EXPLAIN (ANALYZE, FORMAT JSON)} on 100 Q1 queries, we measure the actual number of sub-partitions scanned and pruned.

With application-layer Hilbert key predicates:
\begin{itemize}
\item \textbf{Average pruning rate: 89\%}---5.7 of $\sim$59 sub-partitions scanned per query.
\item \textbf{Average candidate buckets: 7.7}---the Hilbert curve maps a 5\% viewport to fewer than 8 non-contiguous buckets.
\end{itemize}

This confirms that Hilbert's locality preservation is effective: a contiguous 2D region maps to a small set of 1D ranges.

\subsection{Cold-Cache Performance}

Under cold-cache conditions (fresh PostgreSQL connections, no buffer warmup):

\begin{center}
\begin{tabular}{lrr}
\toprule
\textbf{Config} & \textbf{p50 (ms)} & \textbf{p95 (ms)} \\
\midrule
Mono & 486 & 1{,}795 \\
SO & 134 & 278 \\
SPDB & \textbf{53} & \textbf{107} \\
\bottomrule
\end{tabular}
\end{center}

SPDB is \textbf{9.1$\times$ faster than Mono} and \textbf{2.5$\times$ faster than SO} under cold cache (Figure~\ref{fig:cold_warm}).
Partition pruning is a \emph{physical} effect: the planner skips partition files entirely, avoiding disk reads.
This is where SPDB's advantage over SO becomes most pronounced---SO must still scan the entire slide partition, while SPDB reads only the relevant Hilbert buckets.

\begin{figure}[t]
\centering
\includegraphics[width=0.85\columnwidth]{cold_warm_cache}
\caption{Warm vs.\ cold cache Q1 p50 latency. SPDB's cold-cache advantage over SO demonstrates the physical benefit of sub-partition pruning.}
\label{fig:cold_warm}
\end{figure}

\subsection{Hilbert Order Sensitivity}
\label{sec:hilbert_sensitivity}

\begin{table}[t]
\centering
\caption{Hilbert order sensitivity: Q1 viewport latency.}
\label{tab:hilbert}
\begin{tabular}{@{}rrrrl@{}}
\toprule
$p$ & \textbf{Grid} & \textbf{p50 (ms)} & \textbf{p95 (ms)} & \textbf{vs.\ $p{=}8$} \\
\midrule
6 & $4{,}096$ & 52 & 120 & $0.78\times$ \\
\textbf{8} & $\mathbf{65{,}536}$ & \textbf{66} & \textbf{142} & \textbf{---} \\
10 & $1{,}048{,}576$ & 59 & 122 & $0.89\times$ \\
12 & $16{,}777{,}216$ & 64 & 124 & $0.98\times$ \\
\bottomrule
\end{tabular}
\end{table}

All four orders achieve similar performance (52--66\,ms p50), confirming the pruning model's prediction that $p$ is not a sensitive parameter within the moderate range.
$p{=}6$ is marginally fastest because its coarser grid ($64 \times 64$ cells) maps viewport regions to fewer, more contiguous bucket ranges: the resulting SQL predicate has fewer OR clauses, and the planner evaluates fewer candidate partitions.
At finer grids ($p{\ge}10$), a 5\% viewport still spans $\sim\!57$ grid cells per dimension regardless of $p$, so additional resolution provides no pruning benefit while the planner faces more partition boundaries.
The practical implication is that practitioners need not precisely tune $p$: any value in $[6, 10]$ yields near-identical latency.
The default $p{=}8$ provides a safe middle ground that would benefit from finer pruning at smaller viewport fractions ($<$1\%) should the workload shift.

\subsection{Hilbert vs.\ Z-order}

Under controlled comparison (identical partition structure, same number of buckets):

\begin{center}
\begin{tabular}{lrr}
\toprule
\textbf{Curve} & \textbf{p50 (ms)} & \textbf{p95 (ms)} \\
\midrule
Hilbert & \textbf{59} & \textbf{134} \\
Z-order & 82 & 194 \\
\bottomrule
\end{tabular}
\end{center}

Hilbert achieves \textbf{28\% lower} median latency than Z-order, substantially exceeding the 8\% advantage typically reported in the literature~\cite{kamel1994hilbert,moon2001analysis}.
This larger gap arises because our application-layer pruning directly exposes the locality difference: Hilbert's fewer boundary crossings produce fewer candidate buckets per viewport, translating directly to fewer partitions scanned.

\subsection{Viewport Size Sensitivity}

\begin{figure}[t]
\centering
\includegraphics[width=\columnwidth]{viewport_sensitivity}
\caption{Q1 p50 latency vs.\ viewport size. SPDB maintains its advantage at all viewport fractions.}
\label{fig:viewport_sens}
\end{figure}

Figure~\ref{fig:viewport_sens} shows Q1 p50 latency at viewport fractions from 1\% to 20\%.
SPDB maintains a consistent advantage over Mono at all viewport sizes.
At 1\% viewport, SPDB achieves 18\,ms (vs.\ 91\,ms for Mono, a 5$\times$ improvement).
At 20\% viewport, SPDB achieves 192\,ms (vs.\ 609\,ms for Mono, a 3.2$\times$ improvement).
The pruning model predicts this degradation: larger viewports touch more Hilbert buckets, reducing the pruning rate.

\subsection{kNN $k$-Sweep}

\begin{figure}[t]
\centering
\includegraphics[width=0.75\columnwidth]{knn_k_sweep}
\caption{SPDB kNN latency across $k$ values. Latency is nearly constant from $k{=}10$ to $k{=}200$.}
\label{fig:knn_sweep}
\end{figure}

Figure~\ref{fig:knn_sweep} shows SPDB kNN latency remains nearly constant ($\sim$15\,ms p50) across $k \in \{10, 25, 50, 100, 200\}$.
The GiST index's distance-ordered scan handles increasing $k$ without significant overhead, as the additional neighbors are typically in nearby index pages.

\subsection{Concurrency}

\begin{table}[t]
\centering
\caption{Concurrent Q1 throughput (30-sec windows).}
\label{tab:throughput}
\begin{tabular}{@{}lrrr@{}}
\toprule
\textbf{Clients} & \textbf{Mono QPS} & \textbf{SO QPS} & \textbf{SPDB QPS} \\
\midrule
1 & 2.6 & 9.2 & 4.5 \\
4 & 22.6 & 41.2 & 13.8 \\
16 & 47.7 & 65.3 & 15.9 \\
32 & 48.6 & 64.5 & 16.8 \\
64 & 47.3 & 63.3 & 16.3 \\
\bottomrule
\end{tabular}
\end{table}

SO achieves the highest throughput (65\,QPS at 16 clients) because its 29 slide-level partitions provide sufficient isolation without the planning overhead of SPDB's 857 leaf partitions.
SPDB's throughput is constrained by two factors: (1)~each query has dynamically generated \texttt{hilbert\_key} predicates that prevent prepared-statement caching, and (2)~the planner must evaluate 857 partitions per query even when most are pruned at compile time.

Mono saturates at 48\,QPS as all clients contend for the same global GiST index.
SPDB's individual query latency advantage (63\,ms vs.\ Mono's 159\,ms) is offset by per-query planning overhead at high concurrency.

\subsection{Workload Mix}

Under realistic mixed workload (70\% Q1, 15\% Q2, 10\% Q3, 5\% Q4), SPDB sustains \textbf{16.4\,QPS} with per-query p50 latencies of: Q1 = 50\,ms, Q2 = 21\,ms, Q3 = 98\,ms, Q4 = 42\,ms.

\subsection{Storage Overhead}

\begin{table}[t]
\centering
\caption{Storage overhead across storage layouts.}
\label{tab:storage}
\scriptsize
\begin{tabular}{@{}lrrrr@{}}
\toprule
\textbf{Config} & \textbf{Table (GB)} & \textbf{Index (GB)} & \textbf{Total (GB)} & \textbf{B/row} \\
\midrule
Mono & 15.7 & 3.8 & 19.6 & 499 \\
Mono-T & 15.7 & 3.8 & 19.6 & 499 \\
SO & 14.3 & 3.8 & 18.2 & 463 \\
SPDB & 14.7 & 5.0 & 19.8 & 504 \\
SPDB-Z & 14.2 & 5.0 & 19.3 & 492 \\
\bottomrule
\end{tabular}
\end{table}

SPDB's index overhead is $\sim$32\% higher than Mono (5.0\,GB vs.\ 3.8\,GB) due to four per-partition indexes across 857 leaf partitions.
Total storage is comparable ($<$1\% overhead), confirming that the performance gains come at minimal storage cost.
SO achieves the smallest footprint because its 29 partitions reduce per-partition index overhead.

\subsection{Density Analysis}

Object density varies from 105 to 1{,}704 nuclei per megapixel across the 29 slides (Figure~\ref{fig:density}).
This $16\times$ variation in density means that bucket counts range from 5 to 108 per slide.
The Hilbert partitioning naturally adapts to density: denser slides get more buckets, ensuring roughly uniform bucket sizes ($\sim$50K objects) regardless of slide-level density variation.

\begin{figure}[t]
\centering
\includegraphics[width=\columnwidth]{density_distribution}
\caption{Distribution of object density (left) and objects per slide (right) across 29 TCGA BLCA slides.}
\label{fig:density}
\end{figure}

%% ============================================================
\section{Discussion}
\label{sec:discussion}

\textbf{Design trade-offs.}
The five-configuration comparison reveals that performance improvements layer:
(1)~slide-level partitioning provides the largest single improvement ($\sim$2.5$\times$ over Mono);
(2)~Hilbert sub-partitioning adds tail-latency reduction (p95: 141\,ms vs.\ 162\,ms for SO) and dramatically improves cold-cache performance (2.5$\times$ over SO);
(3)~the cost is increased index overhead and reduced concurrent throughput.

\textbf{When to use each configuration:}
\begin{itemize}
\item \textbf{SO} is optimal for high-concurrency multi-user deployments where throughput matters more than individual query latency.
\item \textbf{SPDB} is optimal for single-user or low-concurrency deployments where query latency (especially under cold cache) is the primary concern.
\item \textbf{Mono} is sufficient for small datasets ($<$5M objects) where GiST index depth is manageable.
\end{itemize}

\textbf{Hilbert curve advantage.}
The 28\% Hilbert advantage over Z-order is larger than commonly reported because our pruning mechanism directly translates locality preservation into reduced partition scans.
This validates the theoretical analysis of Hilbert's superior locality properties~\cite{moon2001analysis}.

\textbf{Limitations.}
(1)~Hilbert encoding assumes 2D rectilinear coordinates; 3D or non-Euclidean geometries require extensions.
(2)~Application-layer pruning requires the query layer to compute Hilbert ranges, adding complexity.
(3)~The formal pruning model assumes uniform object density; highly skewed distributions may shift optimal parameters.
(4)~We benchmark on a single node; distributed deployments would need additional evaluation.
(5)~SPDB's concurrent throughput is limited by partition planning overhead; this could be mitigated with prepared statements or PgBouncer.

%% ============================================================
\section{Conclusion}
\label{sec:conclusion}

\textsc{SpatialPathDB} demonstrates that for TCGA-scale pathology workloads, a single PostgreSQL node can serve interactive viewport queries when storage layout respects spatial locality.
On 42.1M nuclei from 29 TCGA BLCA slides, SPDB achieves 63\,ms median viewport latency (2.5$\times$ faster than Mono) and 53\,ms under cold cache (9.1$\times$ faster).
Application-layer Hilbert key pruning eliminates 89\% of sub-partitions before execution.
A formal pruning model accurately predicts optimal parameters without empirical search.
A controlled comparison confirms Hilbert's 28\% advantage over Z-order.

The design space exploration provides practitioners with clear guidance: slide-only partitioning (SO) maximizes throughput; Hilbert sub-partitioning (SPDB) minimizes latency and cold-cache cost.
Source code and benchmark harnesses are available at \url{https://github.com/asishya/SpatialPathDB}.

%% ============================================================
\bibliographystyle{plain}
\begin{thebibliography}{14}

\bibitem{wang2014high}
F.~Wang, T.~Kurc, T.~Oh, P.~Harrington, J.~Chen, et~al.
\newblock Towards building a high performance spatial query system for large scale medical imaging data.
\newblock In {\em Proc. ACM SIGSPATIAL}, pp.~309--318, 2014.

\bibitem{saltz2018spatial}
J.~Saltz, R.~Gupta, L.~Hou, T.~Kurc, P.~Singh, et~al.
\newblock Spatial organization and molecular correlation of tumor-infiltrating lymphocytes using deep learning on pathology images.
\newblock {\em Cell Reports}, 23(1):181--193, 2018.

\bibitem{aji2013hadoop}
A.~Aji, F.~Wang, H.~Vo, R.~Lee, Q.~Liu, X.~Zhang, and J.~Saltz.
\newblock Hadoop-GIS: A high performance spatial data warehousing system over MapReduce.
\newblock {\em Proc. VLDB Endow.}, 6(11):1009--1020, 2013.

\bibitem{guttman1984r}
A.~Guttman.
\newblock R-trees: A dynamic index structure for spatial searching.
\newblock In {\em Proc. ACM SIGMOD}, pp.~47--57, 1984.

\bibitem{kamel1994hilbert}
I.~Kamel and C.~Faloutsos.
\newblock Hilbert R-tree: An improved R-tree using fractals.
\newblock In {\em Proc. VLDB}, pp.~500--509, 1994.

\bibitem{wang2011pais}
F.~Wang, J.~Kong, L.~Cooper, T.~Pan, T.~Kurc, W.~Chen, et~al.
\newblock A data model and database for high-resolution pathology analytical image informatics.
\newblock {\em J. Pathology Informatics}, 2:32, 2011.

\bibitem{eldawy2015spatialhadoop}
A.~Eldawy and M.~Mokbel.
\newblock SpatialHadoop: A MapReduce framework for spatial data.
\newblock In {\em Proc. IEEE ICDE}, pp.~1352--1363, 2015.

\bibitem{sedona2021}
J.~Yu, Z.~Zhang, and M.~Sarwat.
\newblock Spatial data management in Apache Spark: The GeoSpark perspective and beyond.
\newblock {\em GeoInformatica}, 25(1):21--45, 2021.

\bibitem{lawder2000calculation}
J.~K.~Lawder.
\newblock Calculation of mappings between one and $n$-dimensional values using the Hilbert space-filling curve.
\newblock Tech.\ Report JL1/00, Birkbeck College, 2000.

\bibitem{moon2001analysis}
B.~Moon, H.~V.~Jagadish, C.~Faloutsos, and J.~H.~Saltz.
\newblock Analysis of the clustering properties of the Hilbert space-filling curve.
\newblock {\em IEEE Trans.\ Knowl.\ Data Eng.}, 13(1):124--141, 2001.

\bibitem{arge2004optimal}
L.~Arge, M.~de~Berg, H.~Haverkort, and K.~Yi.
\newblock The priority R-tree: A practically efficient and worst-case optimal R-tree.
\newblock In {\em Proc. ACM SIGMOD}, pp.~347--358, 2004.

\bibitem{postgis2023}
PostGIS Development Team.
\newblock PostGIS 3.6 documentation.
\newblock \url{https://postgis.net/docs/}, 2024.

\bibitem{sfc_cost2023}
V.~Pandey, A.~van~Renen, et~al.
\newblock Efficient cost modeling of space-filling curves.
\newblock {\em arXiv:2312.16355}, 2023.

\end{thebibliography}
\end{document}
